% ==============================================================================
% PLANTILLA DE PRESENTACION - UnADM
% Inspirada en la estructura visual de la plantilla beamer de IIIEPE, pero
% adaptada a la identidad academica y al flujo de contenido de los reportes UnADM.
%
% Uso recomendado:
% 1. Duplica este archivo por materia o actividad.
% 2. Actualiza los metadatos editables: titulo, actividad, asignatura, docente,
%    matricula, fecha y bibliografia de apoyo.
% 3. Usa esta presentacion como traduccion visual del reporte UnADM:
%    resumen -> objetivo, introduccion -> problema, desarrollo -> ejes,
%    postura -> lectura critica, conclusion -> cierre, referencias -> fuentes.
% 4. No pegues la planeacion completa en las diapositivas. Convierte la consigna
%    en narrativa, criterios y decisiones de contenido.
% 5. Si la actividad exige producto visual adicional, inserta una diapositiva con
%    esquema, cuadro, mapa o evidencia construida en LaTeX/TikZ o como imagen.
% 6. Compilar desde la raiz del repositorio con:
%    .\scripts\latexmk-build.ps1 .\UnADM\presentacion-unadm.tex
%
% Estructura sugerida para cada actividad:
% [ ] Portada institucional.
% [ ] Objetivo y alcance de la presentacion.
% [ ] Contexto del caso, concepto o problema.
% [ ] Desarrollo en 2-4 ejes claros.
% [ ] Producto visual o evidencia central, si la planeacion lo requiere.
% [ ] Postura personal o lectura critica.
% [ ] Conclusion con aplicacion real.
% [ ] Referencias.
% ==============================================================================

\documentclass[
  spanish,
  aspectratio=169,
  xcolor={dvipsnames,table}
]{beamer}

% Lienzo ampliado 16:9: mantiene la proporcion panoramica y da mas espacio
% util que el tamano por defecto de beamer (12.8 cm x 7.2 cm).
\geometry{paperwidth=19.2cm,paperheight=10.8cm}

% Idioma, tipografia y recursos visuales
\usepackage[utf8]{inputenc}
\usepackage[T1]{fontenc}
\usepackage[spanish,es-tabla]{babel}
\usepackage[scaled=.96]{helvet}
\renewcommand{\familydefault}{\sfdefault}
\usepackage{booktabs}
\usepackage{graphicx}
\usepackage{ragged2e}
\usepackage{tikz}
\usepackage{hyperref}
\usetikzlibrary{calc,positioning,fit,shadows.blur}

% ------------------------------------------------------------------------------
% DATOS EDITABLES
% ------------------------------------------------------------------------------
\newcommand{\studentname}{Martin Jonathan de la Cruz}
\newcommand{\studentshort}{M. J. de la Cruz}
\newcommand{\studentid}{ES2611202040}
\newcommand{\studentemail}{martin.delacruz@unadmexico.mx}
\newcommand{\universityname}{Universidad Abierta y a Distancia de Mexico}
\newcommand{\facultyname}{Licenciatura en Derecho}
\newcommand{\coursename}{Nombre de la asignatura}
\newcommand{\coursecode}{000}
\newcommand{\activitytitle}{Titulo de la presentacion}
\newcommand{\activitysubtitle}{Actividad X - Nombre de la asignatura}
\newcommand{\activityweek}{Semana X}
\newcommand{\teachingfigure}{Nombre de la figura docente}
\newcommand{\deliverydate}{\today}
\newcommand{\locationname}{Roma Norte, Ciudad de Mexico}
\newcommand{\generalbib}{bibliografia-unadm.bib}
\newcommand{\coursebib}{reporte-unadm.tex y bibliografia de la materia}
\newcommand{\departmentlogo}{img/departamentos/UnADM.pdf}
\newcommand{\campusurl}{https://www.unadmexico.mx/}

% ------------------------------------------------------------------------------
% PALETA VISUAL
% Direccion visual: institucional sobria, derivada del isotipo local y pensada
% para conservar compatibilidad con el estilo de los reportes UnADM.
% ------------------------------------------------------------------------------
\definecolor{unadmGreenDark}{HTML}{174A3A}
\definecolor{unadmGreen}{HTML}{5F8F3A}
\definecolor{unadmGold}{HTML}{B88A2A}
\definecolor{unadmPaper}{HTML}{F6F7F2}
\definecolor{unadmInk}{HTML}{1F2A24}
\definecolor{unadmMuted}{HTML}{5D6B66}

\mode<presentation>{
  \usetheme{default}
  \usefonttheme{professionalfonts}
  \setbeamertemplate{navigation symbols}{}
  \setbeamertemplate{blocks}[rounded][shadow=false]
  \setbeamercovered{invisible}
}
\setbeamersize{text margin left=0.55cm,text margin right=0.55cm}

\hypersetup{
  colorlinks=true,
  urlcolor=unadmGreenDark,
  linkcolor=unadmGreenDark
}

\setbeamercolor{background canvas}{bg=white}
\setbeamercolor{normal text}{fg=unadmInk,bg=white}
\setbeamercolor{structure}{fg=unadmGreenDark}
\setbeamercolor{frametitle}{fg=white,bg=unadmGreenDark}
\setbeamercolor{title}{fg=white,bg=unadmGreenDark}
\setbeamercolor{subtitle}{fg=white}
\setbeamercolor{block title}{fg=white,bg=unadmGreen}
\setbeamercolor{block body}{fg=unadmInk,bg=unadmPaper}
\setbeamercolor{alerted text}{fg=unadmGold}
\setbeamerfont{title}{size=\LARGE,series=\bfseries}
\setbeamerfont{frametitle}{size=\large,series=\bfseries}
\setbeamerfont{block title}{series=\bfseries}
\setbeamertemplate{itemize item}{\textcolor{unadmGreen}{\large$\blacktriangleright$}}
\setbeamertemplate{itemize subitem}{\textcolor{unadmGold}{\scriptsize$\blacksquare$}}
\setbeamertemplate{enumerate item}{\textcolor{unadmGreenDark}{\insertenumlabel.}}

\setbeamertemplate{frametitle}{
  \nointerlineskip
  \begin{beamercolorbox}[wd=\textwidth,ht=1.02cm,dp=0.22cm,leftskip=0.35cm,rightskip=0.25cm]{frametitle}
    \begin{minipage}[c]{0.72\textwidth}
      \insertframetitle
    \end{minipage}
    \hfill
    \raisebox{-0.1cm}{\includegraphics[height=0.60cm]{\departmentlogo}}
  \end{beamercolorbox}
  \vspace{-0.04cm}
  {\color{unadmGold}\rule{\textwidth}{1.3pt}}
}

\setbeamertemplate{footline}{
  \leavevmode%
  \hbox{%
    \begin{beamercolorbox}[wd=.34\paperwidth,ht=2.7ex,dp=1ex,leftskip=1em]{author in head/foot}%
      \color{white}\studentshort
    \end{beamercolorbox}%
    \begin{beamercolorbox}[wd=.40\paperwidth,ht=2.7ex,dp=1ex,center]{title in head/foot}%
      \color{white}\coursename
    \end{beamercolorbox}%
    \begin{beamercolorbox}[wd=.26\paperwidth,ht=2.7ex,dp=1ex,rightskip=1em plus 1fil]{date in head/foot}%
      \color{white}\coursecode\hfill\insertframenumber/\inserttotalframenumber
    \end{beamercolorbox}%
  }%
  \vskip0pt%
}
\setbeamercolor{author in head/foot}{bg=unadmGreenDark}
\setbeamercolor{title in head/foot}{bg=unadmGreen}
\setbeamercolor{date in head/foot}{bg=unadmGreenDark}

\newcommand{\accentbar}{%
  \begin{tikzpicture}[remember picture,overlay]
    \path[use as bounding box] (0,0) rectangle (0,0);
    \fill[unadmGreen] (current page.south west) rectangle ([xshift=0.56\paperwidth,yshift=0.08cm]current page.south west);
    \fill[unadmGold] ([xshift=0.56\paperwidth]current page.south west) rectangle ([xshift=\paperwidth,yshift=0.08cm]current page.south west);
  \end{tikzpicture}
}

\newcommand{\sectionframe}[1]{%
  \begin{frame}[plain]
    \begin{tikzpicture}[remember picture,overlay]
      \path[use as bounding box] (0,0) rectangle (0,0);
      \fill[unadmGreenDark] (current page.south west) rectangle (current page.north east);
      \node[opacity=0.09,anchor=east] at ([xshift=-0.45cm]current page.east) {\includegraphics[width=0.43\paperwidth]{\departmentlogo}};
      \fill[unadmGreen] ([xshift=1.05cm,yshift=2.25cm]current page.south west) rectangle ([xshift=1.25cm,yshift=6.40cm]current page.south west);
      \fill[unadmGold] ([xshift=1.48cm,yshift=2.25cm]current page.south west) rectangle ([xshift=1.68cm,yshift=5.25cm]current page.south west);
      \node[
        anchor=west,
        align=left,
        text=white,
        text width=0.54\paperwidth,
        font=\fontsize{26}{31}\selectfont\bfseries
      ] at ([xshift=2.08cm,yshift=0.55cm]current page.west) {#1};
      \node[
        anchor=west,
        align=left,
        text=white!78,
        font=\large
      ] at ([xshift=2.08cm,yshift=-0.55cm]current page.west) {\facultyname};
    \end{tikzpicture}
  \end{frame}
}

\newcommand{\unadmsection}[1]{%
  \section{#1}
  \sectionframe{#1}
}

\title[\coursecode]{\activitytitle}
\subtitle{\activitysubtitle}
\author[\studentshort]{\texorpdfstring{\studentname\\\footnotesize Matricula: \studentid\\\href{mailto:\studentemail}{\studentemail}}{\studentname, Matricula: \studentid}}
\institute[UnADM]{\universityname\\\facultyname\\\coursecode\ \textperiodcentered\ \coursename}
\date[\activityweek]{\locationname\\\activityweek\\\deliverydate}

\begin{document}

% Portada institucional
\begin{frame}[plain]
  \begin{tikzpicture}[remember picture,overlay]
    \path[use as bounding box] (0,0) rectangle (0,0);
    \fill[unadmGreenDark,opacity=0.96] (current page.south west) rectangle ([xshift=0.58\paperwidth]current page.north west);
    \fill[unadmGreen,opacity=0.92] ([xshift=0.58\paperwidth]current page.south west) rectangle (current page.north east);
    \node[anchor=center,opacity=0.10] at (current page.center) {\includegraphics[width=0.72\paperwidth]{\departmentlogo}};
    \node[anchor=north west] at ([xshift=0.58cm,yshift=-0.46cm]current page.north west) {\includegraphics[width=2.55cm]{\departmentlogo}};
    \fill[unadmGold] ([xshift=0.60\paperwidth,yshift=1.18cm]current page.south west) rectangle ([xshift=0.93\paperwidth,yshift=1.30cm]current page.south west);
    \fill[white,opacity=0.15] ([xshift=0.60\paperwidth,yshift=0.92cm]current page.south west) rectangle ([xshift=0.84\paperwidth,yshift=1.02cm]current page.south west);
  \end{tikzpicture}

  \vspace*{1.80cm}
  \begin{minipage}{0.55\paperwidth}
    {\color{white}\fontsize{25}{30}\selectfont\bfseries\raggedright \activitytitle\par}
    \vspace{0.22cm}
    {\color{white!86}\large \activitysubtitle}\\[0.32cm]
    {\color{unadmGold}\rule{0.82\linewidth}{1.3pt}}\\[0.42cm]
    {\color{white}\studentname}\\
    {\color{white!82}\footnotesize Matricula: \studentid}
  \end{minipage}
  \hfill
  \begin{minipage}{0.34\paperwidth}
    \raggedleft
    {\color{white}\Large\bfseries UnADM}\\[0.18cm]
    {\color{white!84}\footnotesize \universityname}\\[0.45cm]
    {\color{white}\small \facultyname}\\
    {\color{white!78}\small \coursename}\\[0.45cm]
    {\color{white!74}\footnotesize \coursecode\ \textperiodcentered\ \activityweek}\\
    {\color{white!74}\footnotesize \locationname}\\
    {\color{white!74}\footnotesize \href{\campusurl}{\campusurl}}
  \end{minipage}
\end{frame}

\begin{frame}{Contenido}
  \tableofcontents
\end{frame}

\unadmsection{Identidad Academica}

\begin{frame}{Ficha de la actividad}
  \begin{columns}[T]
    \begin{column}{0.40\textwidth}
      \begin{center}
        \includegraphics[width=0.44\linewidth]{\departmentlogo}\\[0.25cm]
        {\Large\bfseries\color{unadmGreenDark}UnADM}\\
        {\small\color{unadmMuted}\facultyname}
      \end{center}
    \end{column}
    \begin{column}{0.56\textwidth}
      \begin{block}{Datos generales}
        \textbf{Alumno:} \studentname\\
        \textbf{Matricula:} \studentid\\
        \textbf{Asignatura:} \coursename\\
        \textbf{Codigo:} \coursecode\\
        \textbf{Actividad:} \activitysubtitle
      \end{block}
      \begin{block}{Figura docente}
        \teachingfigure\\
        \deliverydate
      \end{block}
    \end{column}
  \end{columns}
  \accentbar
\end{frame}

\begin{frame}{Contexto academico}
  \begin{columns}[T]
    \begin{column}{0.49\textwidth}
      \begin{block}{Enfoque institucional}
        La presentacion funciona como traduccion visual del reporte academico:
        selecciona el problema, organiza los ejes centrales y muestra una
        postura argumentada sin convertir las diapositivas en un documento.
      \end{block}
    \end{column}
    \begin{column}{0.47\textwidth}
      \begin{block}{Ruta de contenido}
        \begin{itemize}
          \item Objetivo y alcance.
          \item Problema, caso o concepto central.
          \item Desarrollo por ejes.
          \item Postura personal.
          \item Conclusion y referencias.
        \end{itemize}
      \end{block}
    \end{column}
  \end{columns}
  \accentbar
\end{frame}

\unadmsection{Enfoque de Contenido}

\begin{frame}{Correspondencia con la plantilla de informe}
  \begin{block}{Logica recomendada}
    La presentacion no debe repetir el reporte completo. Debe condensar su
    estructura para exponer con claridad: resumen del objetivo, problema central,
    desarrollo por ejes, producto o evidencia, postura y cierre.
  \end{block}
  \begin{columns}[T]
    \begin{column}{0.48\textwidth}
      \begin{block}{Secciones del reporte}
        \begin{itemize}
          \item Resumen / abstract.
          \item Introduccion.
          \item Desarrollo.
          \item Postura personal.
          \item Conclusion.
          \item Bibliografia.
        \end{itemize}
      \end{block}
    \end{column}
    \begin{column}{0.48\textwidth}
      \begin{block}{Traduccion a diapositivas}
        \begin{itemize}
          \item Objetivo y tesis.
          \item Contexto y problema.
          \item Dos o tres ejes de analisis.
          \item Producto visual o caso.
          \item Lectura critica.
          \item Cierre y fuentes.
        \end{itemize}
      \end{block}
    \end{column}
  \end{columns}
  \accentbar
\end{frame}

\begin{frame}{Narrativa sugerida}
  \begin{enumerate}
    \item Abrir con una afirmacion fuerte, no con una definicion blanda.
    \item Traducir la consigna a un problema operativo o pregunta central.
    \item Desarrollar solo los ejes que sostienen la postura.
    \item Usar producto visual cuando la actividad lo pida: cuadro, mapa, esquema, linea de tiempo o tabla.
    \item Cerrar con implicacion practica: que cambia en la lectura del tema.
  \end{enumerate}
  \begin{alertblock}{Regla de contenido}
    Cada diapositiva debe responder una pregunta concreta. Si una diapositiva no
    agrega criterio, debe comprimirse o eliminarse.
  \end{alertblock}
  \accentbar
\end{frame}

\unadmsection{Plantilla Editable}

\begin{frame}{Diapositiva tipo: objetivo y problema}
  \begin{block}{Objetivo}
    \textit{Editable:} Analizar, explicar, comparar, interpretar o evaluar el
    tema solicitado por la actividad.
  \end{block}
  \begin{block}{Problema central}
    \textit{Editable:} El problema no se limita a \ldots, sino a \ldots
  \end{block}
  \begin{block}{Tesis de trabajo}
    \textit{Editable:} La posicion que guia esta presentacion es que \ldots
  \end{block}
  \accentbar
\end{frame}

\begin{frame}{Diapositiva tipo: desarrollo por ejes}
  \begin{columns}[T]
    \begin{column}{0.48\textwidth}
      \begin{block}{Eje 1}
        \textit{Concepto} $\rightarrow$ \textit{funcion} $\rightarrow$ \textit{implicacion}.
      \end{block}
      \begin{block}{Eje 2}
        \textit{Problema} $\rightarrow$ \textit{sistema} $\rightarrow$ \textit{consecuencia}.
      \end{block}
    \end{column}
    \begin{column}{0.48\textwidth}
      \begin{block}{Eje 3 o evidencia}
        \textit{Caso, norma, cuadro comparativo, mapa o ejemplo aplicado.}
      \end{block}
      \begin{block}{Criterio de seleccion}
        No acumular datos. Mantener solo lo que sostiene la conclusion.
      \end{block}
    \end{column}
  \end{columns}
  \accentbar
\end{frame}

\begin{frame}{Diapositiva tipo: producto visual}
  \begin{block}{Cuando la actividad pida evidencia grafica}
    Insertar aqui un cuadro comparativo, esquema, mapa conceptual, red de ideas,
    linea de tiempo o tabla sintesis. Si el producto se construye fuera de LaTeX,
    usar captura nitida y agregar una frase que explique su funcion.
  \end{block}
  \begin{center}
    \fbox{\parbox[c][4.2cm][c]{0.84\linewidth}{\centering Espacio reservado para el producto visual o evidencia central}}
  \end{center}
  \accentbar
\end{frame}

\begin{frame}{Diapositiva tipo: cierre}
  \begin{block}{Postura personal}
    \textit{Editable:} Desde mi perspectiva, el punto mas importante no es
    \ldots, sino \ldots, porque \ldots
  \end{block}
  \begin{block}{Conclusion}
    \textit{Editable:} En sintesis, el valor del tema esta en \ldots y su
    aplicacion real se observa en \ldots
  \end{block}
  \begin{alertblock}{Cierre util}
    No terminar con resumen plano. Terminar con consecuencia, criterio o
    aprendizaje transferible.
  \end{alertblock}
  \accentbar
\end{frame}

\unadmsection{Fuentes y Revision}

\begin{frame}{Fuentes y bibliografia}
  \begin{columns}[T]
    \begin{column}{0.48\textwidth}
      \begin{block}{Bibliografia general UnADM}
        \texttt{\generalbib}
      \end{block}
      \begin{block}{Plantilla de informe base}
        \texttt{reporte-unadm.tex}
      \end{block}
    \end{column}
    \begin{column}{0.48\textwidth}
      \begin{block}{Fuentes de la materia}
        \texttt{\coursebib}
      \end{block}
      \begin{block}{Sitio institucional}
        \href{\campusurl}{\campusurl}
      \end{block}
    \end{column}
  \end{columns}
  \begin{alertblock}{Regla}
    La presentacion no sustituye la bibliografia del reporte. Solo sintetiza las
    fuentes que sostienen la exposicion.
  \end{alertblock}
  \accentbar
\end{frame}

\begin{frame}{Checklist de revision}
  \begin{itemize}
    \item La portada y los datos de actividad estan actualizados.
    \item La narrativa responde a la planeacion sin copiarla literalmente.
    \item Cada diapositiva tiene una funcion clara.
    \item El contenido mantiene enfoque personal, criterio y postura.
    \item El producto visual aparece cuando la actividad lo exige.
    \item La conclusion traduce el tema a una consecuencia practica.
    \item Las fuentes finales coinciden con el reporte y la bibliografia usada.
  \end{itemize}
  \accentbar
\end{frame}

\end{document}
